\chapter{Entanglement Entropy and the Replica Method}
\label{ch:cft-entanglement-replica}
\ConventionsEuclideanCFT

\section{Reduced States and Regulator Dependence}

Let a Lorentzian QFT be given on a Cauchy surface \(\Sigma\), and let
\(A\subset\Sigma\) be an open spatial region with complement
\(A^c=\Sigma\setminus\overline A\).  In a finite lattice regularization with
Hilbert-space factorization
\[
  \mathcal H_\epsilon
  =
  \mathcal H_{A,\epsilon}\otimes\mathcal H_{A^c,\epsilon},
\]
a normalized state \(\rho_\epsilon\) has reduced density matrix
\[
  \rho_{A,\epsilon}
  =
  \operatorname{Tr}_{\mathcal H_{A^c,\epsilon}}\rho_\epsilon .
\]
The von Neumann and Renyi entropies are
\begin{equation}
  S_A(\epsilon)
  =
  -\operatorname{Tr}_{\mathcal H_{A,\epsilon}}
  \rho_{A,\epsilon}\log\rho_{A,\epsilon},
  \qquad
  S_n(A;\epsilon)
  =
  \frac{1}{1-n}
  \log\operatorname{Tr}\rho_{A,\epsilon}^n .
  \label{eq:renyi-von-neumann-definitions}
\end{equation}
The continuum QFT supplies local von Neumann algebras associated to spacetime
domains rather than a canonical tensor product factorization for sharply
separated spatial regions; gauge constraints and edge modes require additional
choices even at finite cutoff.  In this chapter
\(S_A(\epsilon)\) denotes a regulated entropy in a specified regulator and
source chart.  Its power divergences depend on that regulator and source
chart.  The universal logarithmic or finite terms isolated below are
regulator-independent after the allowed local geometric counterterms on the
entangling surface have been fixed.

The modular Hamiltonian of the regulated region is
\[
  K_{A,\epsilon}:=-\log\rho_{A,\epsilon}.
\]
It is defined up to the additive constant that enforces
\(\operatorname{Tr}\rho_{A,\epsilon}=1\).  The first variation of the entropy
around a reference state \(\rho_A\) is the entanglement first law
\begin{equation}
  \delta S_A
  =
  \operatorname{Tr}(\delta\rho_A K_A),
  \qquad
  \operatorname{Tr}\delta\rho_A=0.
  \label{eq:entanglement-first-law}
\end{equation}
Indeed,
\(\delta[-\operatorname{Tr}\rho\log\rho]
=-\operatorname{Tr}\delta\rho(\log\rho+1)\), and the last term vanishes by
normalization.  This identity is finite-dimensional at the regulator level and
then can be used in continuum limits for which both sides have defined
universal parts.

\section{Replica Construction}

Assume that the state is prepared by a Euclidean path integral on a manifold
\(M_1\) cut open along \(A\) at Euclidean time \(0\).  The matrix elements of
\(\rho_A\) are obtained by keeping the fields on the two sides of the cut in
\(A\) open and gluing the fields across \(A^c\).  The trace
\(\operatorname{Tr}\rho_A^n\) glues \(n\) copies cyclically along the cut.  The
result is a manifold \(M_n(A)\), or equivalently a codimension-two replica
defect on the entangling surface \(\partial A\).  With the normalization
\(\operatorname{Tr}\rho_A=1\),
\begin{equation}
  \operatorname{Tr}\rho_A^n
  =
  \frac{Z[M_n(A)]}{Z[M_1]^n}.
  \label{eq:replica-partition-function}
\end{equation}
Let \(W_n(A):=-\log Z[M_n(A)]\).  The von Neumann entropy is obtained by
analytic continuation from positive integer \(n\):
\begin{equation}
  S_A
  =
  -\left.
  \partial_n\log\operatorname{Tr}\rho_A^n
  \right|_{n=1}
  =
  \left.
  (n\partial_n-1)W_n(A)
  \right|_{n=1}.
  \label{eq:replica-entropy-from-effective-action}
\end{equation}
The analytic continuation in \(n\) is an additional assumption.  In examples
it is fixed by locality near the conical defect, reflection positivity, and
growth conditions in \(n\), together with the integer replica sequence.

Near a smooth point of \(\partial A\), the replicated geometry has polar
coordinates
\[
  \dd s^2
  =
  \dd r^2+r^2\dd\theta^2+\dd y^i\dd y^i+\cdots,
  \qquad
  \theta\sim\theta+2\pi n,
\]
where the \(y^i\) are tangent to the entangling surface.  Thus the replica
construction is a precise way to compute the response of the QFT to a conical
excess \(2\pi(n-1)\).  Local geometric counterterms supported on
\(\partial A\) account for the scheme-dependent area-law terms.

\section{One Interval in a Two-Dimensional CFT}

Let the CFT live on the complex plane, and let
\[
  A=[u,v],\qquad L=|v-u|.
\]
The \(n\)-sheeted cover for a single interval is uniformized by
\[
  z
  =
  \left(\frac{w-u}{w-v}\right)^{1/n}.
\]
On the \(z\)-plane the vacuum stress-tensor expectation vanishes.  Under the
map \(z=z(w)\), the stress tensor transforms as
\[
  T(w)
  =
  \left(\frac{\dd z}{\dd w}\right)^2 T(z)
  +
  \frac{c_{2d}}{12}\{z,w\},
  \qquad
  \{z,w\}
  =
  \frac{z'''(w)}{z'(w)}
  -\frac32
  \left(\frac{z''(w)}{z'(w)}\right)^2 .
\]
This gives the one-copy expectation value on the cover
\begin{equation}
  \langle T(w)\rangle_{\mathcal R_n}
  =
  \frac{c_{2d}}{24}
  \left(1-\frac1{n^2}\right)
  \frac{(u-v)^2}{(w-u)^2(w-v)^2}.
  \label{eq:interval-cover-stress-expectation}
\end{equation}
The replicated theory has total stress tensor
\(T_{\rm tot}=\sum_{k=1}^nT_k\).  Comparing
\(n\langle T(w)\rangle_{\mathcal R_n}\) near \(w=u\) with the Ward identity
for a primary twist field \(\sigma_n(u)\) gives
\begin{equation}
  h_n=\bar h_n
  =
  \frac{c_{2d}}{24}
  \left(n-\frac1n\right).
  \label{eq:twist-field-dimension}
\end{equation}
Therefore
\begin{equation}
  \operatorname{Tr}\rho_A^n
  =
  C_n
  \left(\frac{L}{\epsilon}\right)^{
  -\frac{c_{2d}}{6}(n-\frac1n)}
  ,
  \qquad
  C_1=1,
  \label{eq:interval-renyi-trace}
\end{equation}
where \(\epsilon\) is the short-distance cutoff and \(C_n\) is a
cutoff-scheme dependent normalization of the twist fields.  Applying
\eqref{eq:replica-entropy-from-effective-action}, or equivalently
\(-\partial_n\log\operatorname{Tr}\rho_A^n|_{n=1}\), gives
\begin{equation}
  S_A
  =
  \frac{c_{2d}}{3}\log\frac{L}{\epsilon}
  -
  \left.\partial_n\log C_n\right|_{n=1}.
  \label{eq:two-dimensional-interval-entanglement}
\end{equation}
The additive constant is regulator dependent.  The logarithmic coefficient is
universal because it is fixed by the stress-tensor transformation law and the
central charge.

Conformal maps give other standard geometries.  On a circle of circumference
\(L_{\rm cyl}\), an interval of length \(\ell\) has
\[
  S_A
  =
  \frac{c_{2d}}{3}
  \log\left(
    \frac{L_{\rm cyl}}{\pi\epsilon}
    \sin\frac{\pi\ell}{L_{\rm cyl}}
  \right)
  +\hbox{constant},
\]
while at inverse temperature \(\beta\) on the infinite line,
\[
  S_A
  =
  \frac{c_{2d}}{3}
  \log\left(
    \frac{\beta}{\pi\epsilon}
    \sinh\frac{\pi L}{\beta}
  \right)
  +\hbox{constant}.
\]
These formulas follow from the transformation law of the twist two-point
function.

\section{Ball Entanglement in Higher-Dimensional CFT}

Let \(B_R\) be the ball \(r<R\) on the Lorentzian time slice \(t=0\) in the
vacuum of a \(D\)-dimensional CFT.  The domain of dependence of the ball is
conformally mapped to the hyperbolic cylinder by
\begin{equation}
  t
  =
  R\frac{\sinh(\tau/R)}{\cosh u+\cosh(\tau/R)},
  \qquad
  r
  =
  R\frac{\sinh u}{\cosh u+\cosh(\tau/R)},
  \label{eq:causal-diamond-hyperbolic-map}
\end{equation}
where \(u\geq0\) is the radial coordinate on \(\mathbb H^{D-1}\).  After
Euclidean continuation \(\tau=-\ii\tau_E\), smoothness at the entangling
surface requires
\[
  \tau_E\sim\tau_E+2\pi R .
\]
Consequently the reduced vacuum state on the ball is unitarily equivalent,
up to the conformal anomaly in even \(D\), to a thermal state on
\(\mathbb R_\tau\times\mathbb H^{D-1}\) at temperature
\[
  T_{\rm hyp}=\frac{1}{2\pi R}.
\]
The modular Hamiltonian is local:
\begin{equation}
  K_{B_R}
  =
  2\pi
  \int_{r<R}\dd^{D-1}x\,
  \frac{R^2-r^2}{2R}\,
  T_{00}(t=0,\vec x)
  +\hbox{constant}.
  \label{eq:ball-modular-hamiltonian}
\end{equation}
The constant is fixed by \(\operatorname{Tr}\ee^{-K_{B_R}}=1\).

The replica partition function becomes a thermal partition function on the
hyperbolic cylinder:
\begin{equation}
  \operatorname{Tr}\rho_{B_R}^n
  =
  \frac{
    Z_{\mathbb S^1_{2\pi nR}\times\mathbb H^{D-1}}
  }{
    Z_{\mathbb S^1_{2\pi R}\times\mathbb H^{D-1}}^{\,n}
  },
  \label{eq:ball-replica-hyperbolic-thermal}
\end{equation}
and therefore
\begin{equation}
  S(B_R)
  =
  \left.
  \left(1-\beta\frac{\partial}{\partial\beta}\right)
  \log Z_{\mathbb S^1_\beta\times\mathbb H^{D-1}}
  \right|_{\beta=2\pi R}.
  \label{eq:ball-entanglement-thermal-entropy}
\end{equation}
The infinite volume of \(\mathbb H^{D-1}\) reproduces the ultraviolet
divergences near the entangling surface.  The universal term is fixed by the
CFT data:
\begin{equation}
  S_{\rm univ}(B_R)
  =
  \begin{cases}
  (-1)^{D/2-1}\,4a_D\log(R/\epsilon), & D\ \hbox{even},\\[4pt]
  (-1)^{(D-1)/2}\,F_D, & D\ \hbox{odd}.
  \end{cases}
  \label{eq:ball-entanglement-universal-term}
\end{equation}
Here \(a_D\) is the Euler-anomaly coefficient in the normalization for which
\(a_2=c_{2d}/12\) and \(a_4=a_{\rm W}\), while
\[
  F_D=-\log|Z_{\mathcal C}(S^D)|
\]
is the finite sphere free energy after local counterterms have been fixed.
For \(D=2\), \eqref{eq:ball-entanglement-universal-term} reproduces
\((c_{2d}/3)\log(R/\epsilon)\) for an interval whose length is proportional to
\(R\).  For \(D=4\), it gives the sphere-entanglement logarithm
\(-4a_{\rm W}\log(R/\epsilon)\).  For \(D=3\), it gives the universal constant
\(-F_3\), while the renormalized disk entropy used in
\eqref{eq:renormalized-disk-entanglement-f} equals \(F_3\) at a conformal
fixed point.

Equations \eqref{eq:ball-modular-hamiltonian}--\eqref{eq:ball-entanglement-universal-term}
are special to balls and half-spaces in the vacuum.  For a general region the
modular Hamiltonian is a nonlocal operator, and the replica defect is the
primary definition of the entropy.  Shape deformations of the defect couple to
displacement operators and stress-tensor insertions; those structures are
part of defect CFT rather than ordinary local-operator data on flat space.
